\documentclass[11pt]{article}
\usepackage[margin=1in]{geometry}
\usepackage{amsmath,amssymb,booktabs,longtable,array}
\usepackage[hidelinks]{hyperref}
\setlength{\parindent}{0pt}
\setlength{\parskip}{0.5em}
\newcommand{\code}[1]{\texttt{\detokenize{#1}}}
\newcommand{\ind}{\mathbf{1}}
\title{LED Route Guidance: Mathematical Model and Code Implementation}
\author{}
\date{}
\begin{document}
\maketitle

\section{Purpose and interpretation}
The optional LED mode represents directional illuminated strips along the
building's modeled walking network. The strips indicate a shortest-distance
route to an eligible first-floor exterior threshold. People retain the ordinary
behavioral utility model: lighting encourages a direction without replacing
individual preferences by the deterministic student policy.

The implemented influence is a \emph{constant utility bonus}, measured on the
same scale as crowd-manager authority. It is not an acceleration, mechanical
force, or inverse-square attraction. In the existing manager model the
received authority decreases exponentially with manager--guest separation;
LED authority has no such spatial attenuation or memory-decay multiplier.

Continuous strip coverage along the modeled walking network is assumed. The
current model has no separate optical intensity, smoke attenuation, strip
visibility radius, electrical failure, or distance-to-strip detection model.
These are explicit scope assumptions, rather than measured lighting properties.

\section{Shared shortest-route controller}
Use the existing directed graph $G=(V,E)$ with nonnegative edge lengths
$L_e$ in metres. Vertices represent row/aisle attachments, circulation
connections, portal departures and named destination entrances. Stair/access
edges use declared traversal lengths. Rectangular and annular regions use the
same geometric distance rules as crowd-manager recommendations and student
routing. Routes respect floor and region connectivity; the elevator is excluded.

Let $B_L(t)$ be the set of portal IDs unavailable to the LED controller.
For a graph state $v$, define
\[
 D_{B_L(t)}(v)=\min_{\pi:v\leadsto V_{\mathrm{exit}}(B_L(t))}
                                  \sum_{e\in\pi}L_e.
\]
Known unavailable transfer edges and final-threshold seeds are excluded.
The reverse multi-source Dijkstra implementation caches fields for repeated
availability states. Unreachable states have no recommended route.

For a person $i$ at its current position and region, let
$\mathcal A_i^{\mathrm{phys}}(t)$ be the physically permitted local next-target
actions. The controller evaluates
\[
 Q_i^L(a,t)=L_i(a,t)+D_{B_L(t)}(v_a),\qquad
 \widehat a_i^L(t)\in\arg\min_{a\in\mathcal A_i^{\mathrm{phys}}(t)}Q_i^L(a,t).
\]
$L_i(a,t)$ includes distance to and through the proposed first action; for a
portal it includes the actual transfer before downstream distance is added.
Stable exit-ID and action-ID ordering breaks equal-distance ties.

Although the instruction is evaluated at a person's position, the controller
uses one common availability state. The implementation deliberately computes
this recommendation without the person's private blocked-portal list. Private
knowledge cannot redraw a physical strip field differently for each guest.

\section{Information supplied to the person}
The LED network indicates a next local target. If that action is physically
available and its indicated portal sequence contains no closure already known
to the person, it is accepted as an LED instruction. A stale instruction known
to be blocked is ignored; the person can continue with ordinary behavioral
choices instead.

An indicated target is added to the offered action set even if the destination
itself is not visible. The modeling interpretation is following local strips
around a bend, not seeing through a wall. The implementation records the
target location as known. It does not give the person the entire building's
target inventory as student mode does. The full recommended sequence is kept
in diagnostic data, while the behavioral policy uses the next-action cue.

No additional strip memory-strength law is applied. At each normal decision
the person receives the current local cue. The chosen target remains committed
between decisions under the existing scheduler. The default three-second
decision interval and immediate reconsideration on relevant stage changes or
local closure discovery remain in place.

\section{Constant behavioral influence}
Let $U_i^0(a,t)$ be the current ordinary utility before guidance. Its behavioral
component remains
\[
 B_i(a,t)=-w_{Ti}T_i(a,t)+w_{Pi}P_i(a,t)
                         -w_{Ni}N_i(a,t)-w_{Si}C_i(a,t).
\]
The full default utility also retains its local-distance reward, panic term
and residual shock. No familiarity weight is restored and no new person-level
weight is sampled for LEDs.

For adjustable magnitude $\kappa_L\geq0$, the LED contribution is
\[
 \boxed{\displaystyle
 G_i^L(a,t)=\kappa_L\,
      \ind\{\text{valid LED instruction exists}\}
      \ind\{a=\widehat a_i^L(t)\}.}
\]
Thus the indicated action receives exactly $\kappa_L$ utility units and all
other actions receive zero LED bonus. There is no factor $1/d$, $1/d^2$,
$e^{-d/\lambda}$ or age attenuation in this contribution.

With optional manager authority $G_i^M(a,t)$, choice becomes
\[
 a_i^*(t)\in\arg\max_{a\in\mathcal A_i^{\mathrm{offered}}(t)}
       \bigl[U_i^0(a,t)+G_i^M(a,t)+G_i^L(a,t)\bigr].
\]
Manager and LED bonuses add when both nonstudent modes are enabled. Their
directions may agree or differ because managers and lighting can have different
closure knowledge. Existing assigned-manager priority is unchanged.

For an alternative $b\ne\widehat a_i^L$, lighting changes the utility difference
by $\kappa_L$. The person prefers the indicated direction precisely when
\[
 U_i^0(\widehat a_i^L,t)-U_i^0(b,t)
 +G_i^M(\widehat a_i^L,t)-G_i^M(b,t)+\kappa_L\geq0,
\]
subject to the policy's tie handling. A finite bonus can be outweighed by
behavioral preferences or perceived queue costs; compliance is not guaranteed.
Under \code{behavioralOnly=true}, the distance/panic/residual terms are removed
as before, while the behavioral dot product and enabled guidance bonuses remain.

The default is $\kappa_L=12$, matching the default manager's limiting bonus
at zero separation and immediately after receipt. For comparison, the manager
model uses a received factor approximately
\[
 G_i^M(a,t)=\kappa_M\gamma_i e^{-d_{im}/\lambda_M}
            e^{-(t-t_r)/\tau_M}\ind\{a=\widehat a_i^M\},
\]
when its floor/range/occlusion and instruction-validity conditions hold.
The numerical match uses the default $\kappa_M=12$ and $\gamma_i=1$;
changing manager strength does not automatically change LED strength.

Setting $\kappa_L=0$ removes both instruction-based target discovery and the
LED bonus. Displayed strip geometry may remain visible in the GUI, but agents
ignore it. This provides a behavioral zero-strength control.

\section{Adaptive and static lighting}
With \code{adaptive=true}, at each simulation tick the controller uses
\[
 B_L(t)=\{p:\ p\text{ is evacuation-allowed and currently disabled}\}.
\]
Scheduled physical closures are applied before this update. A changed set
invalidates the displayed-network cache and selects the corresponding route
field. This assumes instantaneous portal-availability information for the
lighting controller; no sensor or transmission delay is simulated.

The person receives a changed direction at its next decision. The controller
does not insert every remotely closed portal into the person's knowledge.
Guests retain their ordinary local closure discovery. Physical gates always
prevent passage through a disabled threshold, independent of lighting.

With \code{adaptive=false}, the controller freezes the availability state
captured when the LED option is initialized. Later closures can make the
displayed directions stale. A person who knows that the indicated route is
blocked rejects that cue rather than silently changing the shared strip field.

\section{Physical movement and mode precedence}
LED guidance changes action utilities and next-target information only.
Preferred speeds, reaction delays, social forces, collision avoidance, doorway
service capacities, stair queues and the exterior-threshold completion rule
are retained. There is no new term in the social-force acceleration equation.
The reference graph minimizes distance rather than crowd-dependent travel
time, so lighting does not guarantee the fastest evacuation.

Student mode takes precedence: with \code{studentMode=true}, no LED guidance
object, agent LED records, strip display or bonus is active. The LED checkbox
and strength control are disabled while their preferences are retained for a
later nonstudent reset. Students already use their own direct shortest-distance
selection rule and have no active crowd managers.

For nonstudents, LEDs may be used alone or with explicitly enabled managers.
Enabling lighting does not automatically create staff. The ordinary default
world remains unchanged when LED mode is absent or disabled.

\section{Configuration and controls}
\begin{center}
\begin{tabular}{llll}
\toprule
World key under \code{ledStrips} & Default & Units & Effect\\
\midrule
\code{enabled} & false & Boolean & Activate the optional mode\\
\code{strength} & 12 & Utility & Constant nonnegative bonus\\
\code{adaptive} & true & Boolean & Update controller availability\\
\code{showStrips} & true & Boolean & Draw the reference strip network\\
\bottomrule
\end{tabular}
\end{center}
The separate \code{led_settings.json} example sets \code{enabled=true}.
Strength must be finite and nonnegative. \code{showStrips=false} hides the
visualization without changing agent behavior.

In the GUI, enable \emph{LED route strips (apply/reset)}, set
\emph{LED strength (utility bonus)}, and select \emph{Reset / apply controls}.
Strength is applied on reset rather than changed silently during a run.
The GUI spinner offers 0--100 for convenience; the model validator does not
impose an upper bound on a manually supplied finite nonnegative strength.
\begin{verbatim}
python run.py --n 300 --led-strips --led-strength 12 --gui
python run.py --n 300 --led-settings led_settings.json
python run.py --n 300 --led-strips --led-strength 6
python run.py --n 300 --led-settings led_settings.json --no-led-strips
\end{verbatim}
\code{--led-strength} changes the parameter but does not enable the option
by itself. Student precedence also applies to CLI and world-JSON runs.

\section{Code organization and update sequence}
\begin{longtable}{p{.30\textwidth}p{.64\textwidth}}
\toprule
File & Responsibility\\
\midrule
\endhead
\code{led_guidance.py} & Validates settings; owns controller availability,
instructions, constant bonuses, reference paths and LED summaries.\\
\code{crowd_routes.py} & Shared distance graph, local-action attachment,
shortest-route recommendations and knowledge-keyed route cache.\\
\code{engine.py} & Activates LEDs only in enabled nonstudent runs; updates
the controller after scheduled closures and before decisions; exports state.\\
\code{behavior.py} & Offers the indicated target, applies the LED bonus to
ordinary utility and records selected/counterfactual choices.\\
\code{plan_viewer.py} & Renders green directed strips on the correct floor,
using a revision-dependent visualization cache.\\
\code{app.py} & GUI toggle, strength control, student-mode disabling and
active-mode diagnostics.\\
\code{cli.py}, \code{config.py} & Command-line overrides and validation.\\
\code{led_settings.json} & Separate opt-in settings example.\\
\bottomrule
\end{longtable}

Within \code{LEDGuidance}, \code{before_tick()} updates availability;
\code{instruction()} obtains a valid shared-field direction;
\code{bonus()} returns the constant contribution; and \code{display_paths()}
constructs cached reference strips. At each policy decision,
\code{LocalBehavior.perceived_actions()} adds the cue's target and
\code{LocalBehavior.choose()} evaluates the augmented utilities. The normal
movement engine then advances the chosen action.

\section{Strip visualization and diagnostics}
The display combines shortest graph continuations, reachable radial aisles,
and curved row branches leading to their chosen aisle. Row branches include
the graph's row-to-aisle-centre attachment. Annular routes use straight
segments where valid and tangent/arc paths around the inner disk otherwise.
Arcs are tessellated for drawing, with approximately 0.6-metre maximum arc
sampling spacing. There are no planar lines joining different floors through
stairs. If no exterior route is reachable, the controller displays no
shortest-route strip network.

This is a schematic reference network, not a traced electrical installation.
Actual force-driven trajectories can deviate from its centre paths. Rendering
does not advance the simulation, consume audience RNG draws or alter decisions.

Agent state under \code{data.led} records the current instruction and counts of
recommended, followed and bonus-changed decisions. Each scored alternative
has \code{ledBonus}. The decision also records \code{ledStrength},
\code{choiceWithoutLEDBonus} and \code{ledChangedChoice}. The latter compares
the actual choice with scoring the \emph{same offered actions} after subtracting
only the LED bonus, retaining all other terms, shocks and manager contributions.
It is a local diagnostic, not a full no-lighting causal counterfactual, because
the lighting may also have revealed a target.

Summary exports identify the influence as a constant utility bonus and include
strength, adaptive state, controller closures and adherence counts. Adherence
includes single-action stages and is not an empirically estimated obedience
probability. The mathematical and implementation checks do not calibrate the
default magnitude to observed evacuation behavior. No numerical results are
reported in this write-up.
\end{document}
